\section{定理环境}

\begin{theorem}[线性无关的判定]
向量组 \( \{ \mathbf{v}_1, \mathbf{v}_2, \dots, \mathbf{v}_k \} \) 线性无关当且仅当齐次方程组 \( A\mathbf{x} = \mathbf{0} \) 只有零解，其中 \( A = [\mathbf{v}_1, \mathbf{v}_2, \dots, \mathbf{v}_k] \)。
\end{theorem}

\begin{lemma}[基的扩展]
设 \( \{ \mathbf{v}_1, \dots, \mathbf{v}_k \} \) 是 \( n \) 维线性空间 \( V \) 中的线性无关向量组，则存在 \( n - k \) 个向量将其扩展为 \( V \) 的一组基。
\end{lemma}

\begin{proposition}[维数公式]
设 \( W_1 \) 和 \( W_2 \) 是线性空间 \( V \) 的子空间，则：
\[
\dim(W_1 + W_2) = \dim W_1 + \dim W_2 - \dim(W_1 \cap W_2)
\]
\end{proposition}

\begin{corollary}[直和条件]
子空间 \( W_1 + W_2 \) 是直和当且仅当 \( W_1 \cap W_2 = \{ \mathbf{0} \} \)。此时 \( \dim(W_1 \oplus W_2) = \dim W_1 + \dim W_2 \)。
\end{corollary}

\begin{remark}
在无限维空间中，维数公式不一定成立，因为无限基数的运算规则与有限数不同。
\end{remark}

\begin{example}[子空间示例]
在 \( \mathbb{R}^3 \) 中，xy 平面和 yz 平面的交是 y 轴。因此：
\[
\dim(\text{xy平面} + \text{yz平面}) = 2 + 2 - 1 = 3
\]
\end{example}

\begin{problem}[矩阵求逆]
求矩阵 \( A = \begin{pmatrix} 1 & 2 \\ 3 & 4 \end{pmatrix} \) 的逆矩阵。
\end{problem}

\begin{solution}
\[
A^{-1} = -\frac{1}{2} \begin{pmatrix} 4 & -2 \\ -3 & 1 \end{pmatrix} = \begin{pmatrix} -2 & 1 \\ \frac{3}{2} & -\frac{1}{2} \end{pmatrix}
\]
\end{solution}